\section{Real-data Experimental Protocol}
We evaluate Graph-CoLD on verified real datasets in the formal scope: CICIDS-2017, CESNET-TLS-Year22, UNSW-NB15. CICIDS-2017 uses the full postfilter11 protocol after removing classes below the minimum-count threshold and downsampling the dominant class. CESNET-TLS-Year22 uses a deterministic audit-window subset followed by postfilter25 stratified splitting; we do not describe it as a full-archive evaluation. UNSW-NB15 uses the verified local partition layout with postfilter class policy and temporal plus process/feature-block views.

Each result row records source verification, dataset hash, sample policy, split id, noise seed, model seed, and active views. CICIDS-2017 uses host, IP, and temporal views. CESNET-TLS-Year22 uses IP and temporal views; process and threat-intelligence views are not claimed for this dataset. UNSW-NB15 uses temporal and process/feature-block views because the local partition files do not include source/destination IP columns.

Noise models include clean labels, symmetric label noise, asymmetric label noise, and graph-consistency noise. Noise is injected only into training labels. Graph-consistency rows use the beta values present in the evaluation matrix, and the beta=0 rows are retained to verify consistency with symmetric corruption.

The formal baselines are CoLD, ablation_hard, Noisy-Supervised, Confident-Learning, Co-Teaching, Decoupling, FINE, MCRe, and MORSE. Flash and Argus are not included because they target provenance case-study workflows rather than the formal label-noise matrix. We report Macro-F1, FPR, FNR, tail-class recall, ERR, Tail-ERR, compression ratio, runtime, and memory. Statistical tests are paired by dataset, noise type, noise rate, graph beta, and seed, avoiding unpaired pooled tests.
